\documentclass{article}
\usepackage{amsmath,amssymb,amsthm}
\usepackage{algorithm}
\usepackage[noend]{algpseudocode}
\usepackage{hyperref}
\usepackage{enumitem}
\newtheorem{theorem}{Theorem}
\newtheorem{lemma}{Lemma}
\newcommand{\prox}{\operatorname{prox}}
\newcommand{\R}{\mathbb{R}}

\begin{document}

%%%%%%%%%%%%%%%%%%%%%%%%%%%%%%%%%%%%%%%%%%%%%%%%%
\section*{Proximal Gradient Method (PGM)}
%%%%%%%%%%%%%%%%%%%%%%%%%%%%%%%%%%%%%%%%%%%%%%%%%

\section*{Mathematical Formulation}
Let $f:\R^{d}\!\to\!\R$ be convex and differentiable with $L$‑Lipschitz continuous gradient,
\[
\|\nabla f(x)-\nabla f(y)\|\le L\|x-y\|,\qquad\forall x,y\in\R^{d}.
\]
Let $g:\R^{d}\!\to\!\R\cup\{+\infty\}$ be a proper, closed, convex (possibly non‑smooth) function.
Define the composite objective
\[
F(x)\;:=\;f(x)+g(x),\qquad x\in\R^{d}.
\]

For a stepsize $\eta>0$ the proximal operator of $\eta g$ is
\[
\prox_{\eta g}(v)\;:=\;\arg\min_{u\in\R^{d}}\Big\{g(u)+\frac{1}{2\eta}\|u-v\|^{2}\Big\}.
\]

The **Proximal Gradient update** is
\begin{equation}
\label{eq:update}
x_{k+1}\;=\;\prox_{\eta g}\bigl(x_{k}-\eta\,\nabla f(x_{k})\bigr),\qquad k=0,1,\dots
\end{equation}
The algorithm requires a stepsize $\eta\in(0,2/L)$; the usual choice in theory is $\eta\le 1/L$.

%%%%%%%%%%%%%%%%%%%%%%%%%%%%%%%%%%%%%%%%%%%%%%%%%
\section*{Pseudocode}
%%%%%%%%%%%%%%%%%%%%%%%%%%%%%%%%%%%%%%%%%%%%%%%%%
\begin{algorithm}[H]
\caption{Proximal Gradient Method}
\begin{algorithmic}[1]
\Require Convex $f$, convex $g$, stepsize $\eta\in(0,2/L)$, initial point $x_{0}\in\R^{d}$, max iterations $K$.
\For{$k=0$ \textbf{to} $K-1$}
    \State Compute gradient $g_{k}\gets\nabla f(x_{k})$.
    \State Gradient step: $y_{k}\gets x_{k}-\eta\,g_{k}$.
    \State Proximal step: $x_{k+1}\gets\prox_{\eta g}(y_{k})$.
\EndFor
\State \textbf{return} $x_{K}$.
\end{algorithmic}
\end{algorithm}

%%%%%%%%%%%%%%%%%%%%%%%%%%%%%%%%%%%%%%%%%%%%%%%%%
\section*{Dry Run Example}
%%%%%%%%%%%%%%%%%%%%%%%%%%%%%%%%%%%%%%%%%%%%%%%%%
Consider the scalar problem
\[
f(w)=(w-3)^{2},\qquad g\equiv0,
\]
so that $F(w)=f(w)$.  
Take $w_{0}=0$ and stepsize $\eta=0.1$ (note that $L=2$ for $f$, hence $\eta<1/L$).  
Since $g\equiv0$, $\prox_{\eta g}$ is the identity map.

\begin{center}
\begin{tabular}{c|c|c|c}
Iteration $k$ & $w_{k}$ & $\nabla f(w_{k})$ & $w_{k+1}=w_{k}-\eta\nabla f(w_{k})$\\\hline
0 & $0$ & $2(w_{0}-3)=-6$ & $0-0.1(-6)=0.6$\\
1 & $0.6$ & $2(0.6-3)=-4.8$ & $0.6-0.1(-4.8)=1.08$\\
2 & $1.08$ & $2(1.08-3)=-3.84$ & $1.08-0.1(-3.84)=1.464$\\
\end{tabular}
\end{center}
Corresponding objective values:
\[
F(w_{0})=9,\;F(w_{1})=(0.6-3)^{2}=5.76,\;F(w_{2})=(1.08-3)^{2}=3.6864,\;F(w_{3})=(1.464-3)^{2}=2.362.
\]

The values decrease monotonically, illustrating the method’s descent property.

%%%%%%%%%%%%%%%%%%%%%%%%%%%%%%%%%%%%%%%%%%%%%%%%%
\section*{Assumptions}
%%%%%%%%%%%%%%%%%%%%%%%%%%%%%%%%%%%%%%%%%%%%%%%%%
\begin{enumerate}[label=(A\arabic*)]
\item \textbf{Smoothness.} $f$ is convex and $L$‑smooth:
\[
\|\nabla f(x)-\nabla f(y)\|\le L\|x-y\|,\qquad\forall x,y\in\R^{d}.
\tag{A1}
\]
\item \textbf{Convexity of $g$.} $g$ is proper, closed, convex (no smoothness required). 
\tag{A2}
\item \textbf{Stepsize.} The constant stepsize satisfies
\[
0<\eta\le\frac{1}{L}. \tag{A3}
\]
\item \textbf{Existence of a minimizer.} The set of optimal solutions 
\[
\mathcal{X}^{\star}:=\arg\min_{x\in\R^{d}}F(x)
\]
is non‑empty; pick any $x^{\star}\in\mathcal{X}^{\star}$.
\end{enumerate}

%%%%%%%%%%%%%%%%%%%%%%%%%%%%%%%%%%%%%%%%%%%%%%%%%
\section*{Main Convergence Theorem}
%%%%%%%%%%%%%%%%%%%%%%%%%%%%%%%%%%%%%%%%%%%%%%%%%
\begin{theorem}
\label{thm:rate}
Under assumptions (A1)–(A3), the iterates $\{x_{k}\}$ generated by \eqref{eq:update} satisfy for all $k\ge1$
\[
F(x_{k})-F(x^{\star})\;\le\;\frac{\|x_{0}-x^{\star}\|^{2}}{2\eta\,k}.
\tag{T1}
\]
Consequently $F(x_{k})\to F(x^{\star})$ and the convergence rate is $\mathcal{O}(1/k)$.
\end{theorem}

\begin{theorem}
\label{thm:monotone}
(Descent) For any stepsize $\eta\in(0,2/L)$,
\[
F(x_{k+1})\le F(x_{k})-\frac{1}{2\eta}\|x_{k+1}-x_{k}\|^{2}.
\tag{T2}
\]
\end{theorem}

%%%%%%%%%%%%%%%%%%%%%%%%%%%%%%%%%%%%%%%%%%%%%%%%%
\section*{Theoretical Properties}
%%%%%%%%%%%%%%%%%%%%%%%%%%%%%%%%%%%%%%%%%%%%%%%%%
\begin{itemize}
\item \textbf{Non‑expansiveness of the proximal map.} For any $u,v\in\R^{d}$,
\[
\|\prox_{\eta g}(u)-\prox_{\eta g}(v)\|\le\|u-v\|. \tag{P1}
\]
\item \textbf{Stability w.r.t. stepsize.} The bound (T1) deteriorates linearly with $1/\eta$, hence larger stepsizes (up to $1/L$) accelerate convergence, while $\eta>1/L$ may break the guarantee.
\item \textbf{Relation to Gradient Descent.} When $g\equiv0$, $\prox_{\eta g}$ reduces to the identity and (T1) recovers the classical $O(1/k)$ rate for smooth convex gradient descent.
\item \textbf{Comparison to SGD.} Unlike stochastic methods, PGM is deterministic and obtains the same $O(1/k)$ rate without variance terms, but requires full gradient evaluations.
\end{itemize}

\section*{Discussion}
The proof hinges on two elementary facts: (i) the \emph{descent lemma} for $L$‑smooth functions, and (ii) the \emph{firm non‑expansiveness} of the proximal operator (P1). By coupling these, we obtain a simple telescoping inequality that yields the $1/k$ bound. The result holds for any convex $g$, including indicator functions that enforce constraints (projected gradient descent is a special case).

%%%%%%%%%%%%%%%%%%%%%%%%%%%%%%%%%%%%%%%%%%%%%%%%%
\section*{Preliminaries (Definitions and Lemmas)}
%%%%%%%%%%%%%%%%%%%%%%%%%%%%%%%%%%%%%%%%%%%%%%%%%
\begin{lemma}[Descent Lemma]\label{lem:descent}
If $f$ is $L$‑smooth then for all $x,y\in\R^{d}$,
\[
f(y)\le f(x)+\langle\nabla f(x),y-x\rangle+\frac{L}{2}\|y-x\|^{2}.
\tag{L1}
\]
\end{lemma}
\begin{lemma}[Moreau’s decomposition / Firm non‑expansiveness]\label{lem:prox}
For any convex $g$ and $\eta>0$, the proximal map satisfies for all $u,v$,
\[
\|\prox_{\eta g}(u)-\prox_{\eta g}(v)\|^{2}
\le\langle u-v,\prox_{\eta g}(u)-\prox_{\eta g}(v)\rangle.
\tag{L2}
\]
\end{lemma}

%%%%%%%%%%%%%%%%%%%%%%%%%%%%%%%%%%%%%%%%%%%%%%%%%
\section*{Main Proof (Step‑by‑Step Derivation)}
%%%%%%%%%%%%%%%%%%%%%%%%%%%%%%%%%%%%%%%%%%%%%%%%%
We prove Theorem~\ref{thm:rate}. Throughout the proof we fix an arbitrary optimal point $x^{\star}\in\mathcal{X}^{\star}$.

\noindent\textbf{[Step 1] (Optimality condition of the prox).}  
From the definition of $\prox_{\eta g}$,
\[
x_{k+1}= \prox_{\eta g}\bigl(x_{k}-\eta\nabla f(x_{k})\bigr)
\Longleftrightarrow
\frac{1}{\eta}\bigl(x_{k}-\eta\nabla f(x_{k})-x_{k+1}\bigr)\in\partial g(x_{k+1}),
\]
so that
\[
0\in \nabla f(x_{k})+ \frac{1}{\eta}(x_{k+1}-x_{k})+\partial g(x_{k+1}).
\tag{S1}
\]

\noindent\textbf{[Step 2] (Subgradient inequality for $g$).}  
For any $u\in\partial g(x_{k+1})$ and any $z$,
\[
g(z)\ge g(x_{k+1})+\langle u,\,z-x_{k+1}\rangle.
\tag{S2}
\]
Choose $z=x^{\star}$ and $u:=\nabla f(x_{k})+\frac{1}{\eta}(x_{k+1}-x_{k})$ (the element from (S1)). Then
\[
g(x^{\star})\ge g(x_{k+1})+\Big\langle \nabla f(x_{k})+\frac{1}{\eta}(x_{k+1}-x_{k}),\,x^{\star}-x_{k+1}\Big\rangle.
\tag{S3}
\]

\noindent\textbf{[Step 3] (Combine with smoothness of $f$).}  
Using Lemma~\ref{lem:descent} with $x=x_{k}$, $y=x_{k+1}$,
\[
f(x_{k+1})\le f(x_{k})+\langle\nabla f(x_{k}),x_{k+1}-x_{k}\rangle+\frac{L}{2}\|x_{k+1}-x_{k}\|^{2}.
\tag{S4}
\]

\noindent\textbf{[Step 4] (Add (S3) and (S4)).}  
Adding $g(x_{k+1})$ to both sides of (S4) and then adding (S3) yields
\begin{align*}
F(x_{k+1}) &:= f(x_{k+1})+g(x_{k+1})\\
&\le f(x_{k})+g(x^{\star})+\langle\nabla f(x_{k}),x_{k+1}-x_{k}\rangle
   +\frac{L}{2}\|x_{k+1}-x_{k}\|^{2}\\
&\quad +\Big\langle \nabla f(x_{k})+\frac{1}{\eta}(x_{k+1}-x_{k}),\,x^{\star}-x_{k+1}\Big\rangle .
\end{align*}
Rearranging the inner products,
\begin{align*}
F(x_{k+1})-F(x^{\star})
&\le \langle\nabla f(x_{k}),x_{k+1}-x_{k}\rangle
   +\frac{L}{2}\|x_{k+1}-x_{k}\|^{2}
   +\langle\nabla f(x_{k}),x^{\star}-x_{k+1}\rangle
   +\frac{1}{\eta}\langle x_{k+1}-x_{k},x^{\star}-x_{k+1}\rangle .
\end{align*}
Group the terms containing $\nabla f(x_{k})$:
\[
\langle\nabla f(x_{k}),x_{k+1}-x_{k}+x^{\star}-x_{k+1}\rangle
=\langle\nabla f(x_{k}),x^{\star}-x_{k}\rangle .
\tag{S5}
\]

\noindent\textbf{[Step 5] (Use convexity of $f$).}  
Convexity of $f$ gives
\[
f(x^{\star})\ge f(x_{k})+\langle\nabla f(x_{k}),x^{\star}-x_{k}\rangle,
\]
hence
\[
\langle\nabla f(x_{k}),x^{\star}-x_{k}\rangle\le f(x^{\star})-f(x_{k}) .
\tag{S6}
\]

\noindent\textbf{[Step 6] (Algebraic simplification).}  
Insert (S5)–(S6) into the inequality from Step~4:
\begin{align*}
F(x_{k+1})-F(x^{\star})
&\le \bigl[f(x^{\star})-f(x_{k})\bigr] 
   +\frac{L}{2}\|x_{k+1}-x_{k}\|^{2}
   +\frac{1}{\eta}\langle x_{k+1}-x_{k},x^{\star}-x_{k+1}\rangle .
\end{align*}
Note that
\[
\langle x_{k+1}-x_{k},x^{\star}-x_{k+1}\rangle
= \tfrac12\Big(\|x^{\star}-x_{k}\|^{2}-\|x^{\star}-x_{k+1}\|^{2}-\|x_{k+1}-x_{k}\|^{2}\Big).
\tag{S7}
\]
Plugging (S7) gives
\begin{align*}
F(x_{k+1})-F(x^{\star})
&\le f(x^{\star})-f(x_{k})
   +\frac{L}{2}\|x_{k+1}-x_{k}\|^{2}
   +\frac{1}{2\eta}\Big(\|x^{\star}-x_{k}\|^{2}-\|x^{\star}-x_{k+1}\|^{2}
    -\|x_{k+1}-x_{k}\|^{2}\Big) .
\end{align*}
Cancel the $-\frac{1}{2\eta}\|x_{k+1}-x_{k}\|^{2}$ term with part of the $L$‑term, using (A3) $\eta\le 1/L$:
\[
\frac{L}{2}\|x_{k+1}-x_{k}\|^{2}-\frac{1}{2\eta}\|x_{k+1}-x_{k}\|^{2}
\le 0.
\tag{S8}
\]
Thus
\[
F(x_{k+1})-F(x^{\star})
\le f(x^{\star})-f(x_{k})+\frac{1}{2\eta}\Big(\|x^{\star}-x_{k}\|^{2}-\|x^{\star}-x_{k+1}\|^{2}\Big).
\tag{S9}
\]

\noindent\textbf{[Step 7] (Replace $f(x^{\star})-f(x_{k})$ by $F(x^{\star})-F(x_{k})$).}  
Since $F(x)=f(x)+g(x)$ and $g(x^{\star})\le g(x_{k})$ (by optimality of $x^{\star}$), we have
\[
f(x^{\star})-f(x_{k})\le F(x^{\star})-F(x_{k}).
\tag{S10}
\]
Insert (S10) into (S9):
\[
F(x_{k+1})-F(x^{\star})
\le \bigl[F(x^{\star})-F(x_{k})\bigr]
    +\frac{1}{2\eta}\Big(\|x^{\star}-x_{k}\|^{2}-\|x^{\star}-x_{k+1}\|^{2}\Big).
\tag{S11}
\]

\noindent\textbf{[Step 8] (Rearrange to obtain a telescoping inequality).}  
Move $F(x^{\star})-F(x_{k})$ to the left side:
\[
F(x_{k+1})-F(x_{k})\le
\frac{1}{2\eta}\Big(\|x^{\star}-x_{k}\|^{2}-\|x^{\star}-x_{k+1}\|^{2}\Big).
\tag{S12}
\]

\noindent\textbf{[Step 9] (Summation).}  
Sum (S12) from $k=0$ to $k=K-1$:
\[
\sum_{k=0}^{K-1}\bigl[F(x_{k+1})-F(x_{k})\bigr]
\le \frac{1}{2\eta}\Big(\|x^{\star}-x_{0}\|^{2}-\|x^{\star}-x_{K}\|^{2}\Big)
\le \frac{\|x_{0}-x^{\star}\|^{2}}{2\eta}.
\tag{S13}
\]
The left‑hand side telescopes to $F(x_{K})-F(x_{0})$, yielding
\[
F(x_{K})-F(x^{\star})\le \frac{\|x_{0}-x^{\star}\|^{2}}{2\eta K}.
\tag{S14}
\]
Renaming $K$ as $k$ gives precisely the bound (T1) claimed in Theorem~\ref{thm:rate}. $\square$

\noindent\textbf{Proof of Theorem~\ref{thm:monotone}.}  
Starting from (S12) and discarding the non‑negative term $F(x_{k})-F(x^{\star})$, we obtain
\[
F(x_{k+1})\le F(x_{k})-\frac{1}{2\eta}\|x_{k+1}-x_{k}\|^{2},
\]
which is (T2). $\square$

%%%%%%%%%%%%%%%%%%%%%%%%%%%%%%%%%%%%%%%%%%%%%%%%%
\section*{Rate Derivation (With Constants)}
%%%%%%%%%%%%%%%%%%%%%%%%%%%%%%%%%%%%%%%%%%%%%%%%%
From (T1) we have for any $k\ge1$,
\[
F(x_{k})-F^{\star}\;\le\;\frac{D^{2}}{2\eta\,k},
\qquad D:=\|x_{0}-x^{\star}\|.
\]
If the stepsize is chosen as the largest admissible value $\eta=1/L$, the bound becomes
\[
F(x_{k})-F^{\star}\;\le\;\frac{L D^{2}}{2k}.
\]
Thus the constant factor depends linearly on the smoothness constant $L$ and quadratically on the distance of the initializer to the solution set.

%%%%%%%%%%%%%%%%%%%%%%%%%%%%%%%%%%%%%%%%%%%%%%%%%
\section*{Special Cases}
%%%%%%%%%%%%%%%%%%%%%%%%%%%%%%%%%%%%%%%%%%%%%%%%%
\begin{itemize}
\item \textbf{Projected Gradient Descent.} When $g=\iota_{\mathcal{C}}$ is the indicator of a convex set $\mathcal{C}$, $\prox_{\eta g}$ is the Euclidean projection onto $\mathcal{C}$. The theorem therefore recovers the classical $O(1/k)$ rate for projected gradient descent.
\item \textbf{Lasso (Elastic Net).} For $g(x)=\lambda\|x\|_{1}$, $\prox_{\eta g}$ is the soft‑thresholding operator. The same $O(1/k)$ rate holds for the composite $\ell_{2}$‑loss plus $\ell_{1}$‑regularizer.
\item \textbf{Exact minimizer of $g$.} If $g$ is strongly convex with modulus $\mu>0$, the composite function $F$ becomes $\mu$‑strongly convex, and a simple modification of the above proof yields a linear rate $O\big((1-\eta\mu)^{k}\big)$.
\end{itemize}

\end{document}